\documentclass[11pt]{article}
\usepackage[margin=1in]{geometry}
\usepackage{amsmath,amssymb,booktabs,graphicx}
% The first path works when compiling from docs/; the second from the repo root.
\graphicspath{{../results/student_n3_discussion_note/figures/}{results/student_n3_discussion_note/figures/}}

\title{What the $n=3$ Student experiment taught us about Gibbs and RATTLE}
\author{}
\date{August 2026}

\begin{document}
\maketitle

\paragraph{Purpose.}
This is a discussion note rather than a theorem document.  It records the
sequence of diagnostics that led from acceptance-rate comparisons to a simple
topological explanation of the observed $n=3$ behavior.

\section{The original question}

We initially asked whether constrained Gibbs and RATTLE could be related by
tuning their latent acceptance rates, and whether acceptance or jump size could
explain how often a chain left the region in which the released score root
$\mu^\star$ is the global MLE.  The short answer is no: acceptance and movement
are useful within-method diagnostics, but neither identifies global
communication on this manifold.

\section{Acceptance and ESJD}

The two latent acceptance rates refer to different events.  Gibbs reports the
Metropolis acceptance of a proposed pair move in score coordinates.  Even when
that proposal is rejected, the subsequent exact inverse-branch refresh can
change the latent $x$ values.  RATTLE acceptance instead applies to an entire
projected trajectory.  A near-unit RATTLE acceptance rate therefore says that
the numerical trajectory is stable and usually accepted; it does not say that
the trajectory can reach every part of the target.

Raw $x$-space ESJD was also a poor common scale.  Rare Gibbs resampling events
can replace a central inverse by a far-tail inverse, causing the empirical
second moment to be dominated by a handful of enormous jumps.  Score-space
ESJD, computed after
\[
z_i=\frac{x_i-\mu^\star}{k+(x_i-\mu^\star)^2},
\]
is bounded and much more stable.  At $n=3$, RATTLE's mean score-space SJD was
about $4.1$, $2.6$, and $1.2$ times the Gibbs value for $k=1,2,3$, respectively,
yet a centrally initialized RATTLE chain never left the MLE-valid region.  In
the larger grid cells where RATTLE exits did occur, their score-space jumps
were typically smaller than the overall mean.  Thus score displacement alone
does not explain an exit.

\section{The branch mechanism}

The useful representation is
\[
x\longleftrightarrow (z,b),
\]
where $z$ records the score contributions and $b$ records which inner or outer
inverse of the Student score map is used.  The same $z$ can reconstruct very
different datasets, so global-MLE status depends on both $z$ and the branch
configuration.  In the refined $n=3$ Gibbs runs, 96.0\%, 98.9\%, and 100\% of
MLE exits involved at least one branch change for $k=1,2,3$.  Among the
coordinate-level switches occurring on exits, roughly 90--91\% were
inner-to-outer switches.

\begin{figure}[ht]
\centering
\includegraphics[width=\textwidth]{hexagon_branch_sheet_mle_map.pdf}
\caption{The score vector does not determine MLE status: it must be combined
with the inverse-branch sheet.}
\end{figure}

\section{Why three observations are special}

After normalizing $u_i=2\sqrt{k}z_i$, the constraint $u_1+u_2+u_3=0$ inside
$[-1,1]^3$ is a two-dimensional hexagon.  Accounting for inner/outer labels and
the signs of outer coordinates gives 25 finite chambers.  Finite branch merges
join these chambers into exactly seven connected components: a 13-chamber
central component in which $\mu^\star$ is the global MLE, and six two-chamber
tail components in which it is not.  A branch code such as 111 does not name a
component by itself, because different outer-sign chambers with that same code
belong to different components.

\begin{figure}[ht]
\centering
\includegraphics[width=\textwidth]{chamber_component_connectivity.pdf}
\caption{Exact finite-fold adjacency of the 25 chambers.}
\end{figure}

\section{Why RATTLE is trapped and Gibbs is not}

RATTLE follows a continuous path at finite $x$.  It can change an inverse
branch by passing through a finite turning point,
\[
|z_i|=\frac{1}{2\sqrt{k}}
\quad\Longleftrightarrow\quad
|x_i-\mu^\star|=\sqrt{k}.
\]
It cannot change the sign of an outer coordinate without approaching
\[
z_i\to0,\qquad |x_i|\to\infty.
\]
Those infinite ends separate the seven finite components.  Continuous
finite-position trajectories therefore cannot move from the central component
to a tail component, even though ordinary inner/outer switches remain possible.

The exact-state restart experiment made this visible without relying on a
particular initialization.  Across $k=1,2,3$, 36 RATTLE chains initialized at
exact non-MLE Gibbs states all remained non-MLE, and 12 initialized at exact
MLE-valid Gibbs states all remained MLE-valid.  Each chain retained 7,001
states.  Latent acceptance ranged from 0.9986 to 1.0000, with no projection or
reverse-check failures.

\begin{figure}[ht]
\centering
\includegraphics[width=0.84\textwidth]{rattle_exact_restart_trapping.pdf}
\caption{Exact-state restarts retain their initial connected component.}
\end{figure}

The branch-switch geometry is consistent with the same explanation.  Across
31,402 recorded RATTLE branch-switch events, the median normalized endpoint
distance to a finite merge was 0.00136 and the 95th percentile was 0.0154.
For 65,905 Gibbs events, the corresponding values were 0.129 and 0.734.  Gibbs
does not need to trace a continuous path in $x$: conditional inverse-branch
resampling can jump between finite inverse configurations at fixed score
coordinates.  That nonlocal step supplies the component communication that
continuous RATTLE lacks.

\begin{figure}[ht]
\centering
\includegraphics[width=0.72\textwidth]{branch_switch_merge_gap.pdf}
\caption{RATTLE switches branches at finite merges; Gibbs branch refreshes can
occur far from them.}
\end{figure}

\section{Main empirical numbers}

The following refined comparison used 40,000 retained transitions per method
and $k$.  RATTLE was initialized in the central component.

\begin{center}
\begin{tabular}{ccrrrr}
\toprule
$k$ & method & non-MLE occupancy & exits & exit probability & mean score SJD\\
\midrule
1 & Gibbs  & 75.38\% & 2971 & 30.16\% & 0.00445\\
1 & RATTLE & 0       & 0    & 0       & 0.01837\\
2 & Gibbs  & 43.39\% & 1133 & 5.00\%  & 0.00338\\
2 & RATTLE & 0       & 0    & 0       & 0.00874\\
3 & Gibbs  & 8.81\%  & 150  & 0.41\%  & 0.00457\\
3 & RATTLE & 0       & 0    & 0       & 0.00545\\
\bottomrule
\end{tabular}
\end{center}

These RATTLE zeros are component-conditional results, not estimates of zero
global score-target mass outside the MLE-valid region.

\section{Caveat and conclusion}

Both current algorithms target the score-conditioned distribution.  Neither is
yet an exact sampler for conditioning on the selected/global MLE.  In
particular, RATTLE remaining in the MLE-valid $n=3$ component is a consequence
of reducibility, not evidence that it enforces the global-MLE event.  Detailed
balance can hold within each component while a single continuous chain is not
globally irreducible or ergodic.

The practical conclusion is therefore structural: acceptance tuning and ESJD
cannot repair missing component communication.  Gibbs communicates because its
branch refresh is nonlocal in $x$; continuous RATTLE requires an additional
target-preserving branch-jump kernel before it can explore the complete
score-conditioned manifold.

% Additional ready figures, if useful in a longer version:
% \includegraphics{gibbs_exit_likelihood.pdf}
% \includegraphics{exit_occupancy_grid.pdf}

\end{document}
